\documentclass[11pt, a4paper]{article}

\usepackage{textcomp}
\usepackage{subcaption}
\usepackage{hyperref}
\usepackage[inline]{enumitem}

\bibliographystyle{plain}

\title{A Future-proof and Microarchitecture-agnostic Approach to Collecting CPU Hardware Samples}
\author{
    Marius Albrecht \\
    Chair of Computer Architecture and Parallel Systems \\
    Technical University of Munich \\
    \href{mailto:marius.albrecht@tum.de}{\texttt{marius.albrecht@tum.de}}
}
\date{}

\begin{document}
\maketitle

\section{Motivation}

Quick intro to the topic: HPC\textrightarrow performance\textrightarrow profiling\textrightarrow mitos \\
Explanation of the issue with mitos. \\
Problem statement.

\section{Background}
Maybe split this into "Background" (including related work) and "tools"? I feel like both belongs in the background but the nesting is pretty deep considering e.g. \ref{subsub:mitos} ("Mitos") will be pretty long.
\subsection{Theoretical}
\subsubsection{CPU internals}
As much as is neccesary to explain how PMUs and precise sampling (PEBS and kind of IBS) work and what the differences are between different architectures. \\
This will probably be the basics of processor design in general, OOE, pipelines, caches, memory hierarchy, as well as how PMUs work in general. \\
Lots of terms defined here.
\subsubsection{Precise Event sampling}
Precise event sampling shows up as a umbrella term in the literature. Define this, explain its general position in tee context of pther hardware profiling techniques and their practical uses. \\
Positioning of PEBS, IBS as specific implementations of precise sampling.
\subsubsection{PEBS}
\subsubsection{IBS}
Both kind of obvious. Explain all details, especially what's visible for the programmer. Less focuson IBS (?)
\subsection{tools}
As much as neccesary for the next section.
\subsubsection{perf\_event\_open}
\subsubsection{libpfm}
\subsubsection{Mitos}\label{subsub:mitos}
\subsubsection{Caliper}
\subsection{Related work}
TBA mostly. Id like to see
\begin{enumerate*}
    \item Other tools and possible usages of precise sampling
    \item overhead/accuracy studies, in particular what other factors of the system affect accuracy and overhead
    \item maybe other (possibly more esoteric) implementations. Think POWER, RISCV and maybe ARM
\end{enumerate*}

\section{Implementation}\label{sec:impl}
\subsection{Requirements}
already done. Use in \ref{sec:impl} and \ref{sec:eval}.
\subsection{Tech comparison}
Compare mitos with manual config, mitos with libpfm config and Caliper as "backend" solutions.
\subsection{Implementation}
Rough explanation of the design, no "Design" subsection. \\
Explain the whole implementation.

\section{Evaluation}\label{sec:eval}
Go through requirements. \\
Maybe do some more? This might get kind of short but so be it i guess, I want a toplevel "Evaluation" section.

\section{Conclusion}
TBA
\subsection{future work}
TBA

\end{document}
